\section{下界留下的极值问题}\label{sec:extremal}
第~\ref{sec:core-tenth} 节给出了适用于所有容许集合的下界。要确定锐值，还需要一列合法集合的面积趋近同一个值。我们先把这些构造所用的有限语言写清。

% [U002]
记 $u_q=(\cos q,\sin q)$、$n_q=(-\sin q,\cos q)$，并以 $B_R$ 表示以原点为中心、半径为 $R$ 的闭圆盘。
一个\emph{原始有限体}是形如下式的集合：
\begin{equation}\label{eq:upper-original}
 K=\bigcup_{i=1}^{M}\ \bigcup_{q\in I_i}
 \operatorname{conv}\left(0,
 (b_i+[-1/2,1/2])u_q+h_i n_q\right),
 \qquad \bigcup_{i=1}^{M}I_i=\mathbb {RP}^1,
\end{equation}
其中 $I_i$ 是闭方向区间，$b_i,h_i$ 在每条记录上均为常数。
更换射影方向的提升时，把 $(q,b_i,h_i)$ 换成 $(q+\pi,-b_i,-h_i)$，实际线段保持不变。
在共用端点处保留两侧的记录。
以 $A_*$ 表示这些有限体的面积 $|K|$ 的下确界；$M$ 和外包半径都不设上界。


经典 Cunningham--Schoenberg 构造给出星形问题的族上界
\[
 A_*\le \frac{5-2\sqrt2}{24}\,\pi
\]
\cite{CunninghamSchoenberg1965}。原文通过充分大的奇数指标和任意小的正误差实现这一估计，并没有断言极小值达到；下文将从完整针的源轮廓出发研究更细的面积比较。

\begin{lemma}[开包络中的有限恢复]\label{lem:finite-open-recovery}
设 $S$ 是有限外面积的原点星形集，在每个射影方向上都含有一条完整闭单位线段。对任意开集 $O\supset S$，存在原始有限体 $F\subset O$。特别地，对任意 $\varepsilon>0$，可使 $|F|<|S|^*+\varepsilon$。
\end{lemma}
\begin{proof}
对每个方向 $q$，选出一条完整单位线段 $N_q\subset S$。三角形 $T_q=\operatorname{conv}(0,N_q)$ 是 $O$ 内的紧集，退化时也一样。由旋转在这个紧集上的一致连续性，存在 $0<\delta_q<\pi/4$，使 $|\alpha|\le\delta_q$ 时旋转后的整个三角形仍在 $O$ 内。这里不要求对所有 $q$ 有统一的角宽或半径界。

取 $q$ 的一个实数提升 $t_q$。由 $|\alpha|<\delta_q/2$ 定义的开射影弧覆盖方向圆周，因而有有限子覆盖。用相应的、已经验证安全的较宽闭弧 $|\alpha|\le\delta_q$ 定义 $F$，即取这些旋转三角形的并。每一片都是紧参数积的连续像，相对于旋转坐标系的 $b,h$ 保持常数。因此 $F$ 是紧的原始有限体，接缝记录保留，且每个方向都有一条完整的旋转单位线段留在 $F$ 内。又因 $F\subset O$，有 $|F|\le|O|$。最后由外正则性取 $|O|<|S|^*+\varepsilon$，即得面积结论。
\end{proof}
所需紧性属于每个三角形和方向圆周，而不属于整个 $S$；所以引理也适用于无界的 Borel 代表元及不可测输入。它保证单位覆盖和向上的面积逼近；保留一条指定运动所需的包含恢复将在后文另行证明。

若比较涉及某个固定操作，恢复输入还不够：对恢复后输入施加同一操作所得的输出，也必须满足所需面积比较。第~\ref{sec:shear} 节会同时控制这两端。严格改进一个给定构造，或证明候选族的每个成员都能改进，回答的仍是较小的问题：可得的节省可能随复杂度增加而消失。锐值比较需要对其所作用的近极小序列保留足够的统一控制。

原有限类还可以附加一项要求：在体内选出随方向连续、首尾闭合的完整单位针。第~\ref{subsec:upper-coherent} 节将其面积下确界记为 $A_{\mathrm{coh}}$；这个符号始终指原有限相干子类。紧星体作为更广的测试对象会出现在连续化障碍定理中，但不会改变这一下确界的定义。


连续的方向选择也须与一般的时间参数运动区分。去掉共同星形中心后，即使要求方向参数化连续，单位段的并仍可包含在任意小测度的紧集中~\cite[定理~3.3]{JarvenpaaEtAl2011}。因此，本问题的正面积约束不能仅由选择函数的连续性解释。

